% =============================================================================
% Introduction
% =============================================================================

\section{Introduction}
\label{sec:introduction}

The current paper develops a methodology to quantify countries' exposure to climate change-related disasters, and presents 2 examples of its application for two different countries: Costa Rica and Lao People's Democratic Republic. The methodology is implemented by one joint model called \textbf{Probabilistic Geospatial Risk (PGR) model}.

The methodology focuses on evaluating natural disasters as an important channel through which climate change impacts countries' economies. It employs a set of quantitative and probabilistic methods to estimate and predict the impact of these events.

The approach presented here advances the literature in three key ways. First, it employs a Bayesian modeling framework that enables flexible specification of probability distributions, estimation of the full distribution of disaster occurrences and damages---not just mean values---and rigorous quantification of uncertainty. Second, it incorporates a geospatial approach to connect the geographic patterns of disaster events with global climate trends. Third, it integrates several interrelated models, including time series projections of economic and climate variables, estimates of disaster frequency and severity, and spatial modeling of disaster locations. These models come together to produce probabilistic damage curves, quantifying the range of likely disaster damages in a given location, at a given time while propagating uncertainty between models.

Each step is fully Bayesian, allowing incorporation of prior knowledge held by relevant stakeholders. In addition, the structure of the model allows for scenario-based investigation. For example, it allows the incorporation of different climate forecast scenarios. The result is an end-to-end analysis pipeline that transparently generates, for a given country, region, or city, damage distributions fully characterizing tail risks faced by that region. This comprehensive perspective offers a distinct advantage over existing studies by providing a more detailed and nuanced understanding of the risk of climate-related disasters.

\newpage
